% =====================================================================
%  R. Nielsen, EN FORTOLKNING AF DEN HELLIGE HISTORIE
%  In: For Ide og Virkelighed. Et Tidsskrift, udgivet af R. Nielsen, B. Bjørnson og Rud. Schmidt.  Aargang 1870. Første Bind (Bind 2).
%  Kjøbenhavn: Chr. Steen & Søns Forlag; Louis Kleins Bogtrykkeri.  Bind 2, pp. 289-324.  Antiqua.
%
%  TRANSCRIPTION -- GENERATED; DO NOT EDIT BY HAND.  Rasmus Nielsen's eight essays in the volumes 1870-1873 are built by ONE
%  pipeline (ebook-method/for-idee-1870-73/) into a combined working file, which split7.py cuts into seven transcriptions
%  (Aandsvirkelighed and its sequel, Videre om Aandsvirkelighed, are one: Hans, 2026-10-13).
%  Rebuild from the repo root:   sh ebook-method/for-idee-1870-73/rebuild.sh
%
%  STATUS: COMPLETE (2026-10-13): the volume's title page and the whole essay, Bind 2, pp. 289-324, with its notes.
%
%  EDITION.  The only printing.  Text of record: Nasjonalbiblioteket's scan (nb.no) of the volume, ~/bibliotek/Nielsen, Rasmus/1870-1-for-idee-og-virkelighed.pdf (NB item 440261bfdbcab019d4f313aa5894cf6e) (SCANS.tsv),
%  page images fetched by nbfetch.py, 300 dpi colour, NO text layer; printed page N = PDF page N + 4.  No e-book, no
%  Wikisource text, no reprint, no earlier transcription found (checked 2026-10-13: the repo, the web, the Internet Archive,
%  da.wikisource; HathiTrust and the Google Books API could not be queried that day; Google Books holds at least the 1873
%  volume, nielsen/darwinismen's source -- a second copy, not used).
%
%  ACCURACY (2026-10-13; one collation for all eight texts, 231 pp., ~58,000 words).  The DRAFT is NB's own OCR (ALTO, ABBYY
%  FineReader 12), laid over the page images (combine.py) and turned into LaTeX by draft.py; it was collated word by word
%  with an independent tesseract reading of the same images (collate/check.py, witness 'ocr'): 1556 disagreements.
%  check.py's own rules settled 78.  Candidate rules were first scored BY SCRIPT on a text already read at the image
%  (calibrate.py: Et Synspunkt for Darwinismen, same volume as Verdensanskuelse), then measured on a BLIND PILOT of 152 crops
%  (reader A, 52/52 planted controls right); 17 rules passed (NB's glyph faults: o for ø, D/O/U for Ø, i for l, li for h,
%  l for I, opening » read «, words run together or broken apart; tesseract's: aa read an/ag/å, added accents; the
%  printer's space before a mark, house rule) and settled 772; the rest were read at the image BLIND (562 crops on 94
%  sheets, four readers, 188/188 planted controls right), ~160 decided in the main session from the readers' lines and ~40
%  at 300-dpi zooms.  STRUCTURE read at the image (reader K): every footnote whole (30; collated where they stand, then moved
%  to their marks, then compared word by word with the reading: notecmp.py, 9 differences settled at the image), the
%  headings, the section numerals, the essays' last lines and signatures, the three Greek words, the volume title pages.
%  LETTERSPACING (the print's emphasis, \emph): candidate lines by letter gaps and box width, calibrated on the words the
%  word readers had marked spaced (spaced.py): 361 lines read blind (61/62 roman controls), 25 doubtful extents read at a
%  zoom.  SWEEPS after OPEN 0: words run together that both OCRs share (lower-case + capital), the unknown-word sweep (678
%  words, ~20 suspects, 14 read at the image: 4 OCR faults, 7 misprints kept), --moved, the dropped-word list, the
%  W-shortens list of the blind decisions.
%  SAMPLED PAGE CHECK (pagecheck.py, seed 202610137): one page of each essay or part, pp. 2289, 2400, 3305, 4308, 5033,
%  5208, 6514, 8011 of the working numbering (Bind x 1000 + page), read whole and blind, 1965 words: 0 word errors,
%  1 punctuation (a compound's hyphen lost at a page turn: „Ny-|Holland“) and 1 letterspacing („to be or not to be“)
%  error, both mended.  95 % upper bound for word errors: 1.5 per 1000 words.
%  What this cannot see: a misprint both OCRs normalise into a real word.
%  THE TRANSLATION AS A SECOND READING (2026-10-14): the 15 translators listed 74 doubts; ~50 were read at 300-dpi
%  zooms: 13 faults mended (4 compound hyphens set with a space, 4 specks taken for marks, a lost „I“, a colon read as
%  a semicolon, a comma read as a stop, „à“ read „å“, a space before a mark) and 9 false paragraph breaks at page
%  turns removed (left by the notes moved away; ALTO shows the next page flush); 3 more misprints kept (PRINTED AS IS).
%  The ø rule (Hans 2026-10-08): o-forms that are no word are set with ø; „stode“ (= stood) is a word and stays.
%  PRINTED AS IS Bind 2, p. 303: „Fyldelst“ Fyl-|delst for Fyldest (zoom z5)
%  PRINTED AS IS Bind 2, p. 323: „paayises“ for paavises (zoom z4)
%  PRINTED AS IS Bind 2, p. 321: „Apologotiken“ for Apologetiken (zoom z4)
%  PRINTED AS IS Bind 2, p. 313: „Natnr“ for Natur: a turned u (zoom z4)
%  Not transcribed: running heads, folios, sheet signatures („For Ide og Virkelighed. 1870. I. 20“), library stamps; the
%  large initial of each essay is set as a plain letter.  Quotation marks as printed (» « in the text, sometimes a closing
%  »; „ “ in some notes).  The notes are numbered ¹), ²) per page, as printed.  Greek as printed (textalpha).
%
%  TOOLS (ebook-method/for-idee-1870-73/): nbalto.py, combine.py, alto.py, ocr.py, draft.py, rules.py, calibrate.py,
%  sheets.py + reads.py + READER-BRIEF.md, spaced.py + spaced.tsv, structsheets.py, notecmp.py, unknown.py, zoom.py,
%  pagecheck.py + PAGECHECK-BRIEF.md, patches.py (PATCHES, PRENOTES, NOTES, POST, MISPRINTS), decisions.tsv, split7.py,
%  rebuild.sh.
% =====================================================================
\documentclass[12pt,a4paper,openany]{book}
\usepackage[utf8]{inputenc}
\usepackage[T1]{fontenc}
\usepackage{libertinus}
\usepackage{libertinust1math}
\usepackage{textalpha}
\usepackage[danish]{babel}
\usepackage[protrusion=true,expansion=true]{microtype}
\usepackage{setspace}
\usepackage[margin=1.3in]{geometry}
\usepackage{fancyhdr}
\usepackage{hyperref}
\hypersetup{pdftitle={En Fortolkning af den hellige Historie}, pdfauthor={Rasmus Nielsen}, hidelinks}
\pagestyle{fancy}
\fancyhf{}
\fancyhead[LE,RO]{\thepage}
\fancyhead[C]{\textit{R. Nielsen: En Fortolkning af den hellige Historie (1870)}}
\setstretch{1.3}
\usepackage{marginnote}
\renewcommand{\thefootnote}{\arabic{footnote})}   % the print numbers each page's notes ¹), ²)
\newcommand{\opage}[1]{\marginnote{\footnotesize\textit{[#1]}}}
% the essay's heading: title (and subtitle), centred, a short rule; the print's large initial is not reproduced
\newcommand{\essayhead}[2]{\begin{center}{\large #1}\ifx&#2&\else\\[0.3em]{\normalsize #2}\fi\\[0.5em]\rule{3em}{0.4pt}\end{center}\medskip\noindent}
\newcommand{\signature}[1]{\par\smallskip\hfill #1\hspace*{2em}\par}

\begin{document}

% --- front matter: the title page of Bind 2 (PDF 1 of 1870-1-for-idee-og-virkelighed.pdf), read at the image (reader K, m1);
% short wavy rules are set as plain rules; the library stamps are not transcribed.
\begin{titlepage}
\begin{center}
\vspace*{3em}
{\Huge For Ide og Virkelighed.}\\[1em]
\rule{8em}{0.4pt}\\[1em]
{\large Et Tidsskrift,}\\[0.8em]
{\small udgivet}\\[0.4em]
{\footnotesize af}\\[0.6em]
{\large R. Nielsen,}\\[0.4em]
{\large B. Bjørnson og Rud. Schmidt.}\\[1em]
\rule{8em}{0.4pt}\\[0.6em]
{\large Aargang 1870. Første Bind.}\\[0.6em]
\rule{8em}{0.4pt}\\[4em]
Kjøbenhavn.\\[0.4em]
Chr. Steen \& Søns Forlag.\\[1.5em]
{\small Louis Kleins Bogtrykkeri.}
\end{center}
\end{titlepage}

% ===== TEXT: pp. 289-324 =====
\essayhead{En Fortolkning af den hellige Historie.}{}
% --- p. 289 ---
\setcounter{footnote}{0}%
\opage{289}Naar store Spørgsmaal ere i Bevægelse, og de vaabendygtige
Anskuelser rykke i Marken, er det ikke usædvanligt,
at Felttoget aabnes med en skarp Krigserklæring.
Enhver af de stridende Parter betragter da
sit Standpunkt som en klippefast Borg, en uindtagelig
Fæstning, hvorfra der i forskjellig Retning kan gjøres
Udfald mod den fjendlige Hovedstyrke, medens Standpunktet
selv fra alle Sider er sikret mod ethvert Angreb.
Her er ikke Tale om Underhandlinger, mindelige
Overenskomster eller gjensidige Indrømmelser. Hele
Sandheden er paa den ene Side, hele Usandheden paa
den anden; saa gjælder \emph{to be or not to be}: det, der
ikke kan staae, maa falde. Den verdenskloge Besindighed
ynder ikke den Slags Afgjørelser; den antager ubeseet,
at der er Ret og Uret paa begge Sider. Saa
tysser man paa de Stridende og vil gjerne besværge
Uveiret. Men det er forgjæves, naar Aandedrættet er
tungt og Luften kvalm.

Der gives i Kulturlivet Mellemtilstande, da det næsten
seer ud, som om Menneskeslægten af lutter Klogskab
var i Færd med at blive fjollet, da man forklarer
% --- p. 290 ---
\setcounter{footnote}{0}%
\opage{290}Fænomenerne uden at bekymre sig om deres Væsen og
kan besvare alle Spørgsmaal paa dem nær, der egenlig
angaae Sagen. Aandskræfterne slappes; Fornuften døser,
man skulde troe, at den fik sig en Middagslur, medens
tusinde Tanker af vaagen Forstand sværme i Døgnet
som summende Myg i Heden. Det er i saadanne Mellemtilstande,
at Middelveislæren florerer, thi »Middelveien
er bred« og bliver bestandig bredere. En foregiven
kritisk, men i Virkeligheden ukritisk Sammenhoben
af alskens literære Levninger, en ufordøielig Blanding af
uensartede Bestanddele --- Brudstykker fra gamle, knuste,
sammensunkne Theorier --- ansamles og sætte sig
fast som Dæmning mod alle Principer, som opslemmet
Dynd mod Tankernes friske Strøm, som Grums og Grus,
der tilstopper Ideernes Kildevæld. Under slige Omstændigheder
vil det da være nødvendigt, at det kommer til
Brud, til endelig Afregning, og at Udfordringen er paa
Liv og Død.

Men der kan jo ogsaa føres en Strid paa fredelige
Betingelser. Istedenfor at begynde med en Indskærpelse
af de Punkter, der udelukkende betegne Uenigheden,
kan man vende Fremgangsmaaden om og gjøre Begyndelsen
med at udhæve de Forudsætninger, hvorom der
hersker fuldstændig Enighed. Med Enigheden til Udgangspunkt
vil det saa efterhaanden kunne gjøres indlysende,
hvori Uenigheden bestaaer, hvorfra den stammer,
og hvilken Betydning den egenlig har. I nærværende
Tilfælde, hvor Talen er om en Fortolkning af den hellige
Historie og om den Kritik, der danner Fortolkningens
Grundlag, vælge vi den sidst angivne Fremgangsmaade.
Vi gjøre det saameget hellere, som det Skrift, hvortil
% --- p. 291 ---
\setcounter{footnote}{0}%
\opage{291}vi i denne Afhandling agte at henføre Bedømmelsen af
de evangeliske Beretningers Troværdighed, er en Frugt
af mangeaarig omhyggelig Forskning, et Arbeide af videnskabelig
Værdi, der frembyder væsenlige Enighedspunkter,
nemlig: \emph{Fortolkning} af de synoptiske \emph{Evangelier
af Dr. H. N. Clausen} (Kjøbenhavn 1850).

I den Del af Indledningen, som omhandler »Evangeliernes
historiske Karakter« har den høiærværdige Fortolker,
hvad Kritikens videnskabelige Side angaaer, opstillet
Grundsætninger, som vistnok i deres Almindelighed
maae ansees for uomstødelige. Vi kunne sammenfatte
disse Grundsætninger under følgende tre Synspunkter. 1) \emph{Den gammelkirkelige Antagelse, at man}
»\emph{i Gjenstanden for den evangeliske Beretning,
Kristi jordiske Liv og Virksomhed, havde et
rent historisk Stof, en Række af uomtvistelige
Kjendsgjerninger}«, \emph{er kritisk uholdbar.}
Der maa gjøres Forskjel mellem profan og hellig Historie.
»Det historiske Stof som udfylder den evangeliske Beretning,
kan med Hensyn til indre Beskaffenhed ikke
lignes med noget andet\dots{} Det er det Underfulde, det
Uforklarlige, det Overnaturlige, der her er Beretningens
røde Traad, angiver dens Grundtone; selv der, hvor
Underet ikke træder umiddelbart frem, hvor Begivenhederne
for en Tid føie sig ind i Naturforbindelsens
vante Gang, bevæger dog det Hele sig paa Mysteriets
dunkle, hellige Grund«. Fortolkningen maa altsaa bestemmes
ved Opfattelsen af Underet.

2) \emph{At nægte Underets Mulighed er at fornægte
den hellige Historie.} Imidlertid følger heraf
endnu ingenlunde, at man skulde være tilsinds at ville
% --- p. 292 ---
\setcounter{footnote}{0}%
\opage{292}opgive Besiddelsen af en saadan Historie, »forsaavidt
herved forstaaes en Historie om Verdens viseste Lærer
og helligste Mand, eller Besiddelsen af en hellig Skrift,
hvori denne Mands Lære og hans Livs Historie er bleven
opbevaret«. Men Opgaven bliver, at forklare det
Overnaturlige paa naturlig Maade; og for at løse denne
Opgave, kan man gaae to Veie.

Man kan med den Forudsætning, at Fortællingen
indeholder en virkelig Tildragelse, en sandfærdig Begivenhed,
hvor Alt er gaaet naturlig til, søge at forklare
sig Skinnet af det Overnaturlige dels af Evangelistens
Mangel paa Evne til klar og tydelig Fremstilling, dels
af hans indskrænkede Forestilling om Tingenes Sammenhæng
og hans oprindelige Hang til at see det Naturlige
i et overnaturligt Lys. Saaledes skal f. Ex. Jesu Ord til
Peder ved Anledning af den forlangte Tempelskat (Mtth.
17, 27): »Gak hen til Havet, kast en Krog og tag den
første Fisk, som kommer op; og naar du aabner dens
Mund, skal du finde en Stater; tag denne og giv dem
den for dig og mig«, ifølge Dr. Paulus, være et talende
Bevis for Evangelistens sproglige Ubehjælpsomhed; thi
det, han egenlig vil fortælle, er, at Kristus har sagt til
Peder: »Du er Fisker, udkast Fiskekrogen, og af den
første den bedste Fisk, som du løser af Krogen, vil du
sagtens udbringe den fornødne Penge«. Den historiske
Grund, hvorpaa den evangeliske Beretning maa antages
at hvile, bliver ved denne Behandling mere og mere udhulet,
den fyldige Virkeligheds Kjerne svinder bort, og
kun Skallen bliver tilbage.

Saa kom man da --- »i jævn Skridtgang« --- ind
paa den anden Vei, »den mythiserende Behandling«, hvis
% --- p. 293 ---
\setcounter{footnote}{0}%
\opage{293}betydeligste Repræsentant er Strauss. At den hellige
Historie er en Mythe, et Væv af Myther, vil sige, at
det ikke er virkelige Kjendsgjerninger, men kun religiøse
Forestillinger, vi her have for os; »Tanken arbeider paa
Instinktets Vis. Det er den stærke, endnu uudviklede
Trang til at klare sig, hvad der udfylder Sjælens inderste
Væsen, der sætter de aandelige Kræfter i Bevægelse;
det er de sædelig-religiøse Ideer, der arbeide sig ud af
den aandelige Gjæringstilstand frem til Klarhed og afrundet
Bestemthed, idet de træde frem i legemlig Form
som bestemte Personligheder, i en Mangfoldighed af
Skikkelser, Livsstillinger og Bedrifter«. Den hellige Historie
er med andre Ord et Produkt af ubevidst digtende
Fantasi.

3) \emph{Grundbetingelsen for sand historisk Opfattelse
af den hellige Historie er en troende
Anerkjendelse af Underets Mulighed.} »At en
Historie, der netop som historisk Virkelighed er bleven
Vendepunktet i den religiøse Verdenshistorie, og i den
lange Række af Aarhundreder er gaaet over og ind i
det kristelige Liv, skulde lade sig opløse og udspinde i
en idealistisk-spekulativ Symbolik --- at det personlige
Kristusliv, der i alle sine Hovedtrin af Fornedrelse og
Ophøielse, af Lidelse og Forherligelse er bleven Middelpunktet
for en Virkelighed af den høieste Betydning, for
en Samfundstro og et Samfundsliv med verdensomdannende
Magt, for andagtsfuld Betragtning og Tilbedelse,
for Digtning og Konst, for Anprisning og Efterfølgelse,
paa eengang skulde rives ud af den klare Virkeligheds
Dag og stilles hen i et taaget Dæmringslys, idet det ved
en exempelløs og uforklarlig Illusion skulde være skeet,
% --- p. 294 ---
\setcounter{footnote}{0}%
\opage{294}at en Række af symboliserende Fortællinger --- man veed
ikke, naar eller hvorledes --- kom til at gjælde for
virkelig Historie om en lærende og undergjørende, en
lidende og korsfæstet, en opstanden og himmelfaren Frelser
og Saliggjører, Gjenstand for Bekjendelse og Paakaldelse
i altid større Omraade: dette er en saa uhyre
Paastand, at den kristne Verden, selv dens store ligegyldige
Del ikke fraregnet, maatte bringes til at studse,
til at standse eftertænksomt og gaae i sig selv«. Den
hellige Historie er altsaa Underets Historie, en virkelig
Historie om virkelige Undergjerninger.

Om disse Punkter er her, som sagt, fuldstændig
Enighed; men nu er der i de evangeliske Beretninger
ikke blot Noget, der er Gjenstand for videnskabelig
Prøvelse; her er jo ogsaa Noget, det er Gjenstand for
Tro: Historien er som hellig Historie en Troeshistorie.
Idet nu Forholdet mellem Tro og Viden nærmere skal
bestemmes, begynder Uenigheden.

At Tro og Viden maae, som det hedder, forsones
med hinanden, indrømmes af Alle, der antage, at Religion
og Videnskab lade sig forene i een Bevidsthed.
Det er ikke Forsoningens Mulighed, men dens Betingelser,
hvorom Striden er.

Have Tro og Viden samme Princip, eller er her to
indbyrdes forskjellige, absolut ueensartede Principer? det
er Spørgsmaalet. Den høiærværdige Fortolker antager
det Første, vi det Sidste; et uskatterligt Hjælpemiddel
til Sagens Oplysning have vi i den hellige Historie.

Kan den videnskabelige Anvendelse, Forfatteren har
gjort af de opstillede Grundsætninger, virkelig afgive
Bevis for Troens Princip-Enhed med Viden, da see vi
% --- p. 295 ---
\setcounter{footnote}{0}%
\opage{295}os nødsagede til at opgive vor Synsmaade; men viser
det sig, at den foreliggende Fortolkning netop paa det
Allerkraftigste taler for en principiel Modsætning af Tro
og Viden, da vil jo Fortolkeren selv ikke kunne nægte
en saadan Modsætning uden at komme i Strid med sine
egne Grundsætninger.

Den videnskabelige Kritik skal, hvad ogsaa særlig
indskærpes, lægge an paa »en Udskillelse af overordnede
og underordnede, i historisk Henseende mere eller mindre
hjemlede Partier« i Skriften. »Naar saaledes Evangelierne
indeholde, ved Siden af Beretninger, hvis Udspring
kan forfølges op til det apostoliske Øienvidne, enkelte
Fortællinger, som adskille sig fra den egenlige historiske
Kjerne i Evangelierne ved den fuldkomne Uvished angaaende
de Overleverings-Veie, ad hvilke de ere komne
de evangeliske Forfattere tilhænde (saaledes den evangeliske
For-Historie i de første Kapitler af Matthæus
og Lukas), naar der overfor Beretninger, hvor det Foregaaede
og Fortalte træder frem i beskuelig Klarhed,
staae enkelte andre, der røbe en usikker, tvivlende Opfatning
hos det berettende Vidne selv, og saaledes lade
den fortalte Begivenhed henstaae i mystisk Halvmørke
(dette gjælder f. Ex. om Fristelsen i Ørken, om Forklarelsen
paa Bjerget, om Vandringen paa Søen, om
Opstandelsen af de jerusalemske Døde); naar ved Siden
af mange blot tilsyneladende Uoverensstemmelser, hvor
Skinnet af indbyrdes Modsigelse opløser sig ved nærmere
Betragtning i en høiere eller ringere Grad af Nøiagtighed
og Fuldstændighed, findes enkelte Forskjelligheder,
der ikke lade sig hæve paa nogen tilfredsstillende
Maade (som f. Ex. Tidsangivelsen for det sidste Paaske%
% --- p. 296 ---
\setcounter{footnote}{0}%
\opage{296}maaltid); da kan det ikke være den rette Vei --- fordi
det er en Vei, der ikke fører til noget Hvilepunkt ---
naar der her opstilles et Begreb af Guddommelighedens
Karakter, som ved Magtsprog skal nedslaae Tvivlen og
skræmme Indvendingen tilbage. Thi det er at stille
Troens Sag paa den yderste og farligste Spidse, at
fordre det Unaturlige af den Troende: at han skal holde
Øinene tæt tillukkede, eller at han seende skal give Afkald
paa at kunne see«:

Denne Synsmaade maa i videnskabelig Forstand
ansees for fuldkommen berettiget. Den høiærværdige
Fortolker er, idet han paataler Kritikens Ret, saa langt
fra at stille for strenge Fordringer til Videnskaben, at
vi snarere maa beklage, at han i at gjennemføre disse
Fordringer ikke har været streng nok. Men hvad følger
nu heraf med Hensyn paa Troen?

Det er dog mærkeligt, at netop de Afsnit, som ifølge
ovenstaaende Udtalelse høre til de historisk uvisse, for
ikke at sige uhistoriske, alle paa faa Undtagelser nær
ere optagne blandt dem, der i Kirke-Aarets Løb forelæses
til Menighedens Opbyggelse. Holde vi nu med
Forfatteren fast paa den Forudsætning, at Troes-Sandheden
i den hellige Historie skal klares, begrundes og
bestemmes ved videnskabelig Skriftfortolkning, da er det
let at see, hvad heraf følger. Thi heraf følger: at i
\emph{Forhold til som hine evangeliske Beretningers
Troværdighed svækkes i Videnskaben,} maa
\emph{den ogsaa svækkes i Menigheden.}

Skal Videnskaben virkelig udøve nogen principiel
Indflydelse paa Troen, da maa den jo aabenbart være
i Stand til lidt efter lidt at omdanne Menighedens reli%
% --- p. 297 ---
\setcounter{footnote}{0}%
\opage{297}giøse Bevidsthed, gjøre den opmærksom paa Beskaffenheden
af »det historiske Stof i Evangelierne«, oplade
dens Øie for »en Udskillelse af overordnede og underordnede,
i historisk Henseende mere eller mindre hjemlede
Partier». Saalænge Menigheden vedbliver at høre
paa Bebudelsens Evangelium, paa Evangeliet om Kristi
Fødsel, om de Vise fra Østen, der saae Stjernen o. s. v.,
og vel at mærke, hører uden videnskabelig Kritik, vil
den naturligvis opfatte Begivenhederne ganske som de
fortælles, tillægge det Overnaturlige samme historiske
Virkelighed som det Naturlige, og mene, at det Altsammen
er --- »baade vist og sandt i alle Maader«. Men
efterhaanden som Tilhørerne ved Videnskabens Hjælp
sættes dybere ind i den kritiske Betragtning, vil Troen
og dermed Hørelsen forandre sig. Det vil da underforstaaes,
at disse »fragmentariske Fortællinger« ere af
tvivlsom Oprindelse, at »en mythiserende Indklædning«,
en »udsmykkende Indfatning«, er traadt »udfyldende og
formende» til for at bøde paa »det Sparsomme og Vaklende»
i de historiske Overleveringer. Menigheden vil
da, naar den hører Englenes Lovsang, ikke lade sig
bedaare ved en Forestilling om virkelige Engle, men
nøies med de virkelige Hyrder --- thi for Hyrdernes
Virkelighed kan den videnskabelige Kritik indestaae ---
men hvad Englene angaaer, besinde sig paa »den dramatiske
Benyttelse af de jødiske Forestillinger om Aandeverdenens
Fremtræden paa menneskelig Vis i den menneskelige
Verden, i jødisk Form, selv indtil Navngivning
og Rangorden». Den oplyste Menighed vil i Troen
forstaae, at Underet med Stjernen ikke er noget egenligt
Under, men et mærkværdigt Fænomen, sammensat af
% --- p. 298 ---
\setcounter{footnote}{0}%
\opage{298}et naturligt Himmeltegn og et optisk Bedrag; «thi kan
jeg sige om en Stjerne, at den gaaer foran mig, medens
meget mere jeg gaaer i Retningen henimod den, saa
kan jeg ogsaa sige, at den kommer til at staae der
eller der, medens det meget mere er mig, der kommer
til at staae saaledes, at jeg seer den netop der, saasom
over et Hus«.

Der er fra Troens videnskabelige Standpunkt ofte
klaget over, at den personlige Djævel skal være med i
Historien, og at det er Videnskaben yderst vanskeligt
at faae ham dreven ud af de Troendes Bevidsthed. En
væsenlig Hindring ligger uden Tvivl deri, at Menigheden
hvert Aar paa første Søndag i Fasten hører Evangeliet
(Mtth. 4) om Fristelsen i Ørkenen, og --- naturligvis
uden al Kritik --- opfatter Fortællingen som en virkelig
Begivenhed, en væsenlig Del af den hellige Historie.
Men, naar Videnskaben har faaet sit Alterlys tændt og
det »mystiske Halvmørke«, hvori Begivenheden er indhyllet,
begynder at fordele sig, vil Menigheden vistnok
finde »Vægten af de Vanskeligheder, som trykke den
hele Forestilling om Djævelens personlige Realitet, paa
ganske særegen Maade forstærket ved det Forhold, hvori
denne her træder frem til Jesus. Thi for Djævelens Vedkommende
maa Fremgangsmaaden ved den hele vanskelige
Kamp, det Platte i Ordnen og Indleden af Anfægtelserne
synes ligesaa vanskeligt at forene med den
Kløgt og beregnende Snedighed, der her dog skulde
naae sit Høidepunkt, som, fra Jesu Side, Taaligheden
i at fortsætte disse Forhandlinger, og af den onde Aand
at lade sig føre fra Sted til Sted igjennem Prøvelsens
forskjellige Stadier, vanskeligt forenes med Idealet af
% --- p. 299 ---
\setcounter{footnote}{0}%
\opage{299}den sædelige Renhed, der strax maatte have draget
sig ud af Berøring med det Urene«. Er det først ret
gaaet op for den oplyste Menighed, at Fristelseshistorien,
historisk forstaaet, er ved sin »Plathed« lige prostituerende
for Kristus og for Djævelen, saa faaer Troen vel
at tage tiltakke med den videnskabelige Forklaring, at
det Hele »sandsynligvis grunder sig paa en Meddelelse
af Jesus om Noget, han sjælelig har oplevet i Ørkenens
Ensomhed, hvorved det saa for Resten maa lades
ubestemt, hvorvidt Disciplene have opfattet det Meddelte
som noget indre eller som noget ydre Oplevet«.

Dersom man tager den høiærværdige Fortolkers
Erklæring, at det er den videnskabelige Kritik, der skal
prøve, udfinde og begrunde Evangelieberetningernes historiske
Sandhed, for fuld Alvor, da maa man, hvad
ovennævnte Afsnit angaaer, høilig undre sig over vore
Præster. Thi enten det nu er Evangeliet om Fristelsen
eller om Bebudelsen eller hvilketsomhelst andet Evangelium,
der udgjør Dagens Text, saa gaaer det opbyggelige
Foredrag uden videre ud fra den Forudsætning,
at det forelæste Evangelium er i alle Dele historisk
troværdigt. Præsten tier med Vanskelighederne, forbigaaer
de kritiske Betænkeligheder og lader sig ikke
mærke med Noget. Det er, ifølge Fortolkerens Synsmaade,
en aabenbar Krænkelse af den hellige Sandhed.
En saadan Krænkelse af Sandhedspligten kan ikke undskyldes
med Uvidenhed; den maa da have sin Grund
deri, at der paa Prædikestolen fordres »det Unaturlige»
af Prædikanten: »at han her skal holde Øinene tæt tillukkede, eller at han seende skal give Afkald paa at
kunne see«.

% --- p. 300 ---
\setcounter{footnote}{0}%
\opage{300}Vistnok gjælder den Regel, at man ikke maa give
de Enfoldige Forargelse; men Sandhed er dog altid
Sandhed, og Sandheden maa for Dagen. Er Troesprincipet
virkelig Eet med Vidensprincipet, saa kan der
umulig være een Troes-Sandhed i Videnskaben og en
anden i Menigheden, een paa Kathedret og en anden paa
Prædikestolen.

Den nærmeste Vei til at udrydde Overtroen og befordre
den religiøse Sandhedserkjendelse vilde da være:
af den aarlige Evangeliecyklus at udelukke de Afsnit,
om hvis historiske Troværdighed Videnskaben nærer begrundet
Tvivl. Men naar Talen er om fuldstændig Troværdighed
og historisk Paalidelighed i alle væsenlige
Træk, for hvilke Afsnit af den hellige Historie tør Videnskaben
da videnskabelig indestaae?

At »Forklarelsen paa Bjerget» og «Vandringen paa
Søen» ligesaa lidt kunne komme i Betragtning som
»Fristelseshistorien« eller de »fragmentariske Fortællinger«,
der afgive »Forhistorien«, har Fortolkeren selv
allerede antydet.

\emph{Forklarelsen paa Bjerget} (Mtth. 17, Mark. 9,
Luk. 9) \emph{gaaer ud;} det er ikke et Billede med tydelige
Træk af historisk Virkelighed, vi her have for os, men
»et Taagebillede i mystisk Halvlys». Han (Kristus)
\emph{blev forvandlet for Disciplene,} men hvori Forvandlingen
egenlig har bestaaet, kan Kritiken ikke udfinde.
«Hvorledes den antydede Metamorfose er at tænke:
som strækkende sig til Legemets indre Natur --- og i
dette Tilfælde, hvorledes forenelig med Jesu sande Menneskelighed
--- eller som blot tilsyneladende med Hensyn
til det ydre Udseende; --- hvorledes Tilstedevæ%
% --- p. 301 ---
\setcounter{footnote}{0}%
\opage{301}relsen skal tænkes af de tvende Jødedommens Heroer»
(Moses og Elias) »Repræsentanterne for Loven og Profeterne:
virkelig eller visionær; --- hvad Formaal der
skal tillægges den beskrevne Sammenkomst, og hvad der
skal tænkes som Samtalens Gjenstand, --- hvorledes
Røsten fra Skyen og de udtalte Ord skal forklares: alt
dette er Spørgsmaal, paa hvilke et nogenlunde tilfredsstillende
Svar maatte være at finde, for at den fortalte
Begivenhed skulde kunne blive Gjenstand for Forestilling
og Fremstilling. Men af Fortællingens Beskaffenhed
er det indlysende, at og hvorfor intet saadant Svar
kan gives. Thi de trende Øienvidner« (Petrus, Jakobus
og Johannes) »vare betyngede af Søvn, de vidste
ikke, hvad de talede, thi de vare forfærdede«; da de
ret kom til Besindelse og »saae sig om, saae de Ingen
uden Jesus alene«.

\emph{Vandringen paa Søen} (Mtth. 14, Mark. 6)
\emph{gaaer ud.} Fortolkeren fastholder med filologisk Ret,
at her ikke er Tale om en Vandring \emph{ved} Søen eller
\emph{henimod} Søen, men ordret om en Vandring \emph{paa} Søen,
\emph{hen over Søen,} og henviser til Job 9, 8, hvor der
siges om Gud, at han vandrer paa Havet som paa et
Gulv. Hele Fortællingen, siger han, har »øiensynligt
denne Forklaring til Forudsætning; thi medens Disciplene
paa Skibet vare midt ude paa Søen, hedder det, at
Jesus vilde gaae dem forbi, og de forbausedes, menende,
at det var et Spøgelse, de saae for sig«. Forsaavidt
gjør Kritiken da, under Forudsætning af Underets Mulighed,
ingen videnskabelig Indvending mod Begivenheden;
men --- naturam furca --- nu kommer det. Fortællingen
om Peder, der ogsaa vil prøve at vandre paa
% --- p. 302 ---
\setcounter{footnote}{0}%
\opage{302}Vandet, et karakteristisk Optrin, som fuldender Billedet
og giver det en rig symbolsk Betydning, kan Kritiken
ikke godkjende. Det opbyggelige Foredrag taler imidlertid
i en ganske anden Tone. »Dette er skrevet, at
I skulle troe, at Jesus er den Kristus, den Guds Søn;
dette Ord ville vi anvende paa vor forelæste hellige
Fortælling om Herrens Vandring paa Genesareth Sø\dots{}
Naar vi troe paa ham, som har opvakt Døde, har hersket
over Dødens Lov, skulle vi da ikke ogsaa troe paa
ham, som har hersket over Tyngdens Lov? Vi skulle
da ogsaa lære dette af vort Evangelium, at hvad det
kommer an paa, det er at have Fasthed, Bestandighed
i Troen, at vi ikke, ligesom Petrus, i Troen skulle begynde
Vandringen, men naar vi see det stærke Veir,
blive forfærdede, tvivle og synke«. Det er jo meget opbyggeligt.
Men naar Troens og Videns Princip-Enhed
fastholdes, bliver Grundlaget for denne opbyggelige Tale
hult og svævende. Prædikanten taler om Noget, der aldrig
er skeet. Kritiken har »grundet Formodning om«,
at »Episoden om Peder« er en senere Tilsætning, der
skyldes den græske Oversættelse eller Bearbeidelse af
Matthæus-Evangeliet --- »en Formodning, som bestyrkes
ved den hele Karakter af denne Fortælling, der minder
om den udsmykkende og forstærkende Tendens i de apokryfiske
Evangelier«.

»\emph{Stormens Stillen paa den galilæiske Sø}»
(Mtth. 8, Mark. 4, Luk. 8) \emph{gaaer ud, for at Menigheden}
ikke skal tage Forargelse af den paafaldende Misforstaaelse,
der har indsneget sig i Fremstillingen, en
Misforstaaelse, der er saa stor, at den videnskabelige
Fortolkning fuldstændig omdanner Begivenhedens histo%
% --- p. 303 ---
\setcounter{footnote}{0}%
\opage{303}riske Karakter. Alle tre Evangelier lægge nemlig Eftertrykket
eenstemmig paa, at Storm og Bølger maae
lyde Kristi almægtige Ord. Men Videnskaben forstaaer
Sagen anderledes: hele Betydningen ligger i Modsætningen
af Mesterens »ophøiede Ro» og Disciplenes »ængstelige
Forsagthed«, saa at Jesu truende Ord »snarere
vel har været henvendt til Disciplene end til Elementerne,
til Veir og Vind«. --- Der er, naar vi følge den
videnskabelige Kritik, endnu mange Afsnit, som paa
Grund af historisk Upaalidelighed maae udelukkes, ja
saamange, at man ikke kan undertrykke det Spørgsmaal,
om der tilsidst bliver et eneste tilbage.

Det er, vi gjentage det, ingenlunde vor Mening, at
den høiærværdige Fortolker i sine Undersøgelser er gaaet
de historiske Kjendsgjerninger for nær. Tvertimod! naar
Videnskabens Fordring skal i nogen Maade skee Fyldelst,   % PRINTED AS IS (Bind 2, p. 303): „Fyldelst“ Fyl-|delst for Fyldest (zoom z5)
vil vist Ingen kunne behandle den evangeliske
Historie med mere Varsomhed eller større Beredvillighed
til overalt, hvor det er muligt, at lade det Underfulde
træde frem med Præg af historisk Virkelighed,
end skeet er i det her omhandlede Skrift. Men netop
denne Stræben: \emph{ud fra Videnskaben at sikre
og begrunde det Overnaturlige i Troens Indhold},
har, som vi strax skulle see, bevirket en Svækkelse
af Grundsætningerne, en Famlen og Vaklen i Bedømmelsen
af Fænomenerne, der nærmer sig til det
Principløse.

For at imødegaae de Indvendinger, den vantro Kritik
(Strauss, Weisse, B. Bauer) har reist imod Helbredelsen
af Høvidsmandens Tjener (Mtth. 8, Luk. 7) som en
»Helbreden i det Fjerne, uden at der er Tale om den
% --- p. 304 ---
\setcounter{footnote}{0}%
\opage{304}Syges egen Tro«, bemærker Forfatteren: »uagtet det
maa indrømmes, at Tanken har lettest ved at tilegne
sig og fastholde Underet, hvor det fremtræder knyttet
til Troen som Mellemled\dots{}, saa fremtræder dog den
samme evangeliske Under-Typus overalt, hvor den hellige
Villie aabenbarer sig som den mægtige Villie, i den
legemlige Verden som i den sædelige«. Her gjøres altsaa
Aarsagen, den hellige, mægtige Villie, til Miraklets
Kjendemærke.

Men dette Kjendemærke udviskes igjen, hvor Talen
er om Helbredelsen af den blodsottige Kvinde (Mtth. 9,
Mark. 5, Luk. 8). Saaledes som Fortællingen er givet
hos Markus og Lukas, »fører den nemlig med Nødvendighed
til Forestillingen om Jesu undergjørende Kraft
som virkende\dots{} uden at Virksomhedens Grund bliver
at søge i Villiens bestemmende Magt\dots{} Den helbredende
Kraft er altsaa her fremstillet som sat i Bevægelse
ved den blotte udvortes Berøring, og den paafulgte
Helbredelse fremgaaer som et fysisk Fænomen,
hvor uvilkaarlig Tanken paatrænger sig om det elektriske
Apparat, opfyldt af den stærke Kraft, der kun
behøver at nærmes hen imod et Legeme for at strømme
ud til synlige Virkninger«. Strengt taget, burde nu et
saadant Mirakel ligefrem forkastes; men man kan jo i
det Sted forkaste Kjendemærket, opgive den mægtige
Villies bestemte Hensigt, holde sig til den Troende og
gjøre Troen til Under-Kraftens elektriske Leder.

Denne videnskabelige Opfattelse har opbyggelig Hjemmel
i Schleiermachers Prædikener. »Dybt og sindrigt«,
siger Fortolk., »har Schleiermacher yttret sig om dette
Problem: at ligesom den aandelige Indvirken, der ud%
% --- p. 305 ---
\setcounter{footnote}{0}%
\opage{305}gaaer fra den Enkelte, ikkun for en ringe Del har sin
Grund i den bestemte Villie og Hensigt, men et saadant
Omløb finder Sted af de aandelige Kræfter, at det
Bedste kan virkes, uden at den Virkende veed deraf at
sige, saaledes forholdt det sig ogsaa med Jesu underfulde
Gjerninger i det Legemlige, som just derved fjernede
sig fra Tingenes sædvanlige Gang og nærmede sig
til det aandelige Livs Enemærker; hist saavelsom her
var Troen, Attraaen og Længselen den eneste Betingelse
for at blive delagtig i den underfulde Virkning«.

Vil Læseren ikke tage til Takke med den Schleiermacherske
Analogi, men finde Begivenheden urimelig,
saa har den kritiske Fortolkning endnu den naturlige
Udvei, at søge Grunden til det Urimelige i en Misforstaaelse
hos Evangelisterne. Matthæus giver Fortællingen
kort og lægger Helbredelsens Kraft i Jesu Ord:
\emph{Datter, vær frimodig, din Tro haver frelst}
dig. Man kan nu antage, at dette netop er det Historiske,
medens Fortællingen hos Markus og Lukas »gjøre
Indtrykket af traditionelle Udmalinger, dannede ved Reflexion
over Helbredelsens Art og Maade«.

Den naturlige Udvei er imidlertid den farligste. Vel
tager det sig ved første Øiekast ret rimeligt ud, at en
Evangelist i al Troskyldighed kan være kommen til at
benytte et og andet Sagn, hvori det Mirakuløse ved Tilsætninger
og »traditionelle Udmalinger« er blevet misforstaaet:
det er da ikke Evangelisten, men Sagnet, der
er Skyld i Misforstaaelsen. Men see vi nærmere til,
saa forholder Sagen sig anderledes. Thi den Evangelist,
som kan optage Sagn, der stille Miraklet i et falsk
Lys, lægger jo netop derved for Dagen, at han mangler
% --- p. 306 ---
\setcounter{footnote}{0}%
\opage{306}Blik for Underets Særkjende, at han, rentud sagt, ikke
er i Stand til at skjelne mellem sande og falske Mirakler.

At en Evangelist kan have feilet i Angivelsen af
en eller anden historisk Yderside, et Navn, en Tidsbestemmelse
o. desl., er uvæsenligt, thi en Evangelist er
ingen videnskabelig kritiserende Historieforsker; heller ikke
er der just noget Urimeligt i, at han ved at afbenytte Efterretninger
af sagnmæssig, uvis Oprindelse kan i sin Fremstilling
have forvexlet »det indvortes med det udvortes
Oplevede«, thi i religiøs Forstand faaer det Udvortes
altid sin Betydning ved det Indvortes. Men at antage
om en Evangelist, at han i tankeløs Godtroenhed har
godkjendt Fortællinger, der i Virkeligheden forvrænge
Miraklet, gaaer paa ingen Maade an, thi det er at tilintetgjøre
hans Troværdighed som Evangelist. Hvad
Fortolkeren i denne Henseende anfører til Undskyldning
for det lidet Troværdige i »Udmalingen af Underet med
den blodsottige Kvinde, med Vandringen paa Søen, med
Stateren i Fiskens Mund, med Opstandelsen af »de jerusalemske
Døde o.s.v., er ikke Undskyldning, men Anklage;
det er lutter Tvivl om de hellige Forfatteres Evne
til at behandle det historiske Stof, tilslørede Indsigelser
mod deres Troværdighed.

»Vi skulle«, siger den høiærværdige Forf., »af Skriften
selv lade os veilede til det sande Begreb om det
hellige Aabenbaringsord«. Ja, det skulle vi. Men denne
Grundsætning bliver jo glemt eller skudt tilside, hver
Gang vi fordre en ganske anden Fremstilling af det
Vidunderlige, end den, der vitterlig er os given; hver
Gang vi henvise til »senere, skæve, mistydende Reflexioner«
og staae i Begreb med at sætte vor egen saakaldte vi%
% --- p. 307 ---
\setcounter{footnote}{0}%
\opage{307}denskabelige Mirakelforestilling høiere end den evangeliske.

Hvorlidt man, hvad Underet angaaer, kan være tilsinds
at lade sig af Skriften »veilede til det sande Begreb«,
derpaa afgiver Kritiken over Bespisningsmiraklet
(Mark. 6, Mtth. 14, Luk. 9) et talende Exempel. »For
den exegetiske Betragtning har«, siger Forf., »det Under:
at fem Tusinde ere blevne mættede ved fem Brød og to
Fiske, saa fuldstændig udvortes Hjemmel som nogen af
de evangeliske Undergjerninger«. Ikke desto mindre
maa den videnskabelige Kritik føle sig i høi Grad utilfredsstillet
ved dette Under. Grunden er, »at det, som
her træder os imøde, er en Magthandling i absolut Ubegribelighed
af den skabende Guddomskraft\dots{}. et Vidunder
(τέρας) staaer for os, som tillukker sig for enhver
Tankeanstrængelse, afviser ethvert Spørgsmaal, der søger
Besvarelse, som ubesvarligt«. --- Denne usikkre, med sig
selv uenige Opfattelse af det Overnaturlige, denne svævende
Ræsonneren for og imod, disse tvivlraadige Overslag
over rimelige, lidet rimelige og urimelige Begivenheder,
denne ufrugtbare Afveien af sandsynlige, lidet
sandsynlige og usandsynlige Mirakler, er i lige Grad
uheldbringende for Religion og Videnskab: den gjør
Troen tvivlende og Viden principløs.

Det kan af vort betragtede Fortolkningsarbeide
bevises, at den hellige Historie principielt maa udelukke
Videnskaben. Ved at indrømme Underets Mulighed har
Viden sat sin absolute Grændse og hermed sagt sit sidste
Ord. Historien har imidlertid ikke en almindelig Mulighed,
men virkelige Kjendsgjerninger, disse Tegn, disse
Syner, disse Aabenbarelser, disse i bestemte Forestil%
% --- p. 308 ---
\setcounter{footnote}{0}%
\opage{308}linger opfattede Undergjerninger, til Indhold. Skulde
Videnskaben paa nogen Maade være i Stand til at fælde
en begrundet Dom over de enkelte, i bestemte Fænomener
fremtrædende Undere, da maatte den jo først
opdage de Love, hvorefter det Overnaturlige kan yttre
sig og gribe ind i det Naturlige. Er den nu, som vi
\emph{alle vide, ganske uvidende om disse Love, saa er dens}
Kritik uden Maalestok: enten maa den lade Undergjerningerne
staae for, hvad de ere: usandsynlige, ubegribelige,
eller den maa søge at opløse enhver enkelt bestemt
Forestilling om ethvert enkelt bestemt Under i »det
Underfuldes« ubestemte, ubestemmelige Mulighed. I
første Tilfælde vil Videnskaben drage sig tilbage fra den
hellige Historie, i sidste Tilfælde vil dens Kritik opløse
den hellige Historie.

Vor høiærværdige Fortolker har maaske ikke besindet
sig paa, at det virkelige Resultat af hans for
Videnskaben betydningsfulde Arbeide er en fuldbyrdet
Opløsning af den evangeliske Historie. Og dog er Opløsningen
uimodsigelig. Der hersker i denne Fortolkning
fra Først til Sidst en saadan Skilsmisse mellem det
Naturlige og det Overnaturlige, at kun det Naturlige
bliver historisk; medens det Overnaturlige --- det i sig
Ubestemte og Ubestemmelige --- unddrager sig enhver
mulig historisk Opfattelse. Det hedder i Indledningen
til første Afsnit, at den videnskabelige Kritik vil »betrygge
den historiske Kjerne« ved at skille den fra »den
udsmykkende Indfatning« og navnlig »hævde den historiske
Virkelighed af et underfuldt Grund - Faktum ved
Kristi Indtræden i Verden«; men den, som stoler paa,
at han her ved Kritikens Bistand skal finde enten Be%
% --- p. 309 ---
\setcounter{footnote}{0}%
\opage{309}viser for det underfulde Grund-Faktums Virkelighed eller
nærmere Oplysninger om dets historiske Beskaffenhed, bliver
høiligen skuffet.

\emph{Hvorledes er det egenlig gaaet til med Bebudelsen
af Johannes Døberens Fødsel?} Man
veed det ikke. Lukas (1, 5—25) beskriver det rigtignok
udførlig; ved at anføre de nærmere Omstændigheder og
Bi-Omstændigheder mener Evangelisten vel, at han klarer
os Forestillingen om det i Begivenheden Overnaturlige og
fortæller en sandfærdig Historie. Men Fortolkeren er af
en anden Mening. Han mener, at det, Evangelisten her
fortæller, ikke er nogen virkelig Begivenhed, ikke nogen
sandfærdig Historie, men en gammeltestamentlig Udsmykning
af et Familiesagn. »Ligesom det er sandsynligt,
at traditionelle Sagn ere benyttede, som opbevaredes
i Johannes’s Familie, saaledes er det øiensynligt, at det
Gl. Test.s Fortællinger om Simsons og om Samuels
Fødsel (Domm. 13. 1 Sam. 1) have tjent til Forbillede
i de enkelte Træk«. Hvad er det saa, der egenlig er
skeet? Det veed man ikke: Begivenheden selv falder
udenfor Historien.

\emph{Hvorledes er det gaaet til med Bebudelsen
af Kristi Fødsel?} Det veed man ikke. Matthæus
(1, 18—25) og Lukas (1, 26—38) mene jo rigtignok at
kunne meddele tilforladelig Oplysning og give os hver
sit Bidrag; men det er en from Indbildning. Kritiken
beviser, at der ingen Sammenhæng er imellem »de tvende
Beretninger«, eftersom den ene aabenbart gjør den anden
overflødig. »Intet er vel i sig selv naturligere, end
at der om en Begivenhed, hvor det Overnaturlige tabte
sig i hemmelighedsfuld Dunkelhed, vare forskjellige For%
% --- p. 310 ---
\setcounter{footnote}{0}%
\opage{310}tællinger i Omløb; og i alle Tilfælde har der, fra Beretternes
Side, ikke været tænkt paa at forene dem
indbyrdes; thi det lader sig ikke nægte, at --- efter
Engleaabenbarelsen for Maria --- Drømmesynet for Josef
staaer lidet motiveret«.

\emph{Kan Englenes Aabenbarelse for Hyrderne}
(Luk. 2, 8—20) \emph{ansees for en historisk Begivenhed}?
Ja, paa Prædikestolen, men ikke i Videnskaben!
Vel er her et Grund-Faktum; men naar Kritiken har
faaet »Indfatningen« skilt fra »Kjernen«, saa forsvinder
det vidunderlige Grund-Faktum i Sagnets Dunkelhed og
taber sig tilsidst i en sindrig psykologisk Tvetydighed.
Vi maae nemlig med Schleiermacher »antage, at denne
Historie, først efter at Jesu Rygte er bleven begrundet,
er ved Enkeltes Erindring bleven fremdraget af Forglemmelsen;
men at den hverken er en af Intet sammenpustet
Luftblære, eller et forsætligt Opspind, derfor borger
os den Lethed, som de Kristne i den Egn havde ved at
udspørge Maria, som nødvendigen maatte vide Besked
dermed, eller de Disciple, med hvilke hun mest havde
levet, om Sagens Beskaffenhed«. Altsaa Noget maa der
være skeet, noget Underfuldt, men hvad er det for Noget?
Det veed man ikke. Begivenhedens Øiemed er
nemlig ikke at søge i »den fjernere Virkning, som den
frembragte, saa at sige efter dens Opstandelse, men
kun i den umiddelbare paa Hyrderne selv. Privataabenbarelsen
for Hyrderne kan nu forklares paa en ganske
naturlig og dog næsten overnaturlig Maade deraf, at En
eller Anden iblandt dem har, »i sin særegne Fromhed«,
udbedet sig et Tegn. »Dette blev dem, som Historien
lærer, til Del, og i denne, i saa store Tider ikke sjeldne,
% --- p. 311 ---
\setcounter{footnote}{0}%
\opage{311}Tilfredsstillelse ogsaa af Enkeltes Længsel ligger det
Mærkværdige og Guddommelige«.

\emph{I sit offenlige Liv er Kristus selv det underfulde
Grund-Faktum;} men at Kritiken heller
ikke her kan finde sig tilrette med de historiske Fænomener
af dette Grund-Faktum, have vi allerede seet i
mange Exempler. For ikke at forkaste alle enkelte
Undergjerninger, men dog holde til Raade med Mangfoldigheden,
tilkjender og frakjender den ræsonnerende
Forstand Frelseren naturlige og overnaturlige Gaver, som
det bedst kan falde sig. At Jesus, efter hvad vi have
hørt om Høvidsmandens Tjener, vel har havt Gaver til
at helbrede en Syg i det Fjerne, bliver stiltiende indrømmet;
men at han ogsaa har havt Gaver til at kunne
see et Æsel i det Fjerne, bliver ikke indrømmet. »Æselet«
--- hedder det om Indtoget i Jerusalem (Mtth. 21,
Mark. 11, Luk. 9) --- »afhentes af de tvende Disciple
i den lige overfor beliggende Landsby, sandsynlig Bethfage«.
Alle tre Evangelister ere jo vistnok enige om,
at lægge Vægt paa den underfulde Forudsigelse af, at
de skulde »finde« Æselet, og saaledes antyde en »mysteriøs
Sammenhæng ved den hele Tildragelse«. Men heri
kan Kritiken ikke give de hellige Forfattere Medhold.
Videnskaben maa finde det sandsynligt, at der »enten af
Jesus er bleven tilføiet nærmere fornøden Oplysning,
eller ogsaa en forud truffen Aftale fra Jesu Side maa
underforstaaes«.

Endnu et Træk. Fortolkeren lader sig forlyde med,
at Forklarelsen paa Bjerget vanskelig kan forenes med
»Jesu sande Menneskelighed«; her kan altsaa ikke tænkes
paa noget overnaturligt Grund-Faktum; men kan
% --- p. 312 ---
\setcounter{footnote}{0}%
\opage{312}\emph{Opstandelsen og Himmelfarten} da bedre forenes
med »Jesu sande Menneskelighed«? Her maae vi holde
fast paa et mysteriøst Faktum. Men det mysteriøse
Faktum er saa langt fra at klare sig i Fremstillingen,
at det tvertimod gjør Fremstillingen uklar. »De indbyrdes
Afvigelser, med hvilke de fire evangeliske Beretninger
om Opstandelsen gjennemkrydse hverandre«, virke
i den Grad forstyrrende, at »der kun med stor Usikkerhed
lader sig, ved den indbyrdes Sammenstilling, tilveiebringe
en nogenlunde sammenhængende Fortælling«;
med andre Ord: det, der egenlig skulde fortælles, Tildragelsernes
overnaturlige \emph{Hvorledes}, bliver ikke fortalt;
det, som her skulde bevidnes, bliver ikke bevidnet.
Og det er ganske i sin Orden, det kan ikke være anderledes.
Grunden til Usikkerheden ligger ingenlunde i
nogen Uagtsomhed eller Forsømmelse fra Øienvidnernes
Side, den ligger i Begivenhedens Natur, i »det overmægtige
Indtryk, hvoraf de første Øienvidner maatte gribes,
hvorved det navnlig kommer i Betragtning, at Traditionen
i sine enkelte Træk her har havt sin første Hjemmel i
den usikkre Iagttagelse af de kvindelige Vidner«.

Som det gaaer med Opstandelsens Evangelium, saaledes
ogsaa med \emph{Himmelfartens.} Himmelfarten bevidnes
som Grund-Faktum, som Artikel i en Troesbekjendelse,
men om Himmelfarten selv gives der ingen
Historie. Og hvorfor ikke? Fordi det vidunderlige
Grund-Faktum staaer i aabenbar Strid med al Historie.
Hvad der ikke kan fremtræde i en bestemt Form paa
en bestemt Maade, kan heller ikke indtræde paa et bestemt
Sted til en bestemt Tid, kan ikke blive Begivenhed;
men hvad der efter sin Natur ikke kan blive Begivenhed,
% --- p. 313 ---
\setcounter{footnote}{0}%
\opage{313}kan heller ikke blive historisk. Kritiken sætter en
uoverstigelig Kløft imellem Himmelfartens Faktum og
»den \emph{Form og Maade,} hvorpaa dette Faktum, denne
underfulde Indgaaen i den himmelske Herlighed, skal
tænkes fuldbyrdet«. Ja, Himmelfarten er, naar vi see
nærmere til, ikke engang noget egenligt Faktum; det er
et »tillukket Mysterium«, men et tillukket Mysterium er
intet Faktum. »Havde der« --- hermed slutter den høiærværdige
Fortolker --- »ikke her været et tillukket
Mysterium allerede for Apostlene, da vilde hverken Forbigaaelsen
af det største σημεῖον« (Tegn) »i de tvende
apostoliske Evangelier være forklarlig, eller den Dunkelhed,
hvori Hentydningen i de tvende andre Evangelier
er holdt«. Dette er Enden paa den hellige Historie
--- i Videnskaben.

Dette Mærkværdige, at en saa klar, saa skarpsindig,
i kritisk Tænkning saa øvet, med sit store videnskabelige
Apparat saa fortrolig Forfatter, som Dr. H.
N. Clausen, er ved at gjennemføre sine Grundsætninger
kommen til et Resultat, der netop er det Modsatte af
det, hvorpaa han oprindelig har sigtet, vækker Formodning
om, at der i Grundsætningerne selv maa skjule sig
en bedaarende Uklarhed, en blændende Tvetydighed.
Vi skulle i Korthed bestræbe os for at afsløre denne
Tvetydighed.

At der, ifølge Grundsætningen (1), maa gjøres Forskjel
mellem profan og hellig Historie, er uimodsigeligt;
men Spørgsmaalet er, hvori Forskjellen bestaaer og hvor
dybt den grunder i Sagen. Den »historiske Troværdighed«
er af en anden Natnr, af en anden Oprindelse og har en   % PRINTED AS IS (Bind 2, p. 313): „Natnr“ for Natur: a turned u (zoom z4)
ganske anden Grund i den hellige Historie end i den profane.

% --- p. 314 ---
\setcounter{footnote}{0}%
\opage{314}Til Begrebet \emph{Troværdighed} svarer Begrebet Tro;
det Troværdige er Troens Gjenstand. Heraf følger, at
Troværdighedsbegrebet maa forandres, naar Troesbegrebet
forandres. Men nu er der paa det evangelisk-protestantiske
Dannelsestrin dog vel neppe Nogen, der vil
benægte, at den religiøse Tro (fides religiosa) er væsenlig
forskjellig fra hvad man i almindelig Forstand
kalder historisk Tro (fides historica).

Den blot historiske Tro angaaer Begivenheder, der
i sin Tid have været Gjenstand for udvortes, umiddelbar
Erfaring, og om hvilke Øienvidnet forsaavidt kan
antages at have havt en tilforladelig Viden. Begivenheder
af denne Natur høre alle ind under den lovbestemte
Tingenes Orden, der danner de udadgaaende Iagttagelsers,
de sandselige Erfaringers faste Grundlag. Ved
Opfattelsen af slige Begivenheder vil det Ydre overalt
være bestemmende for det Indre. Øienvidnet har at
holde sig til, hvad der har viist sig, gjengive det Oplevede,
som et Speil giver Gjenstandens Billede, ikke
indblande Noget af sit Eget, ikke sætte Farve paa
Kjendsgjerningerne. Ved at opfylde disse Fordringer
er et Øienvidne troværdigt. Den, der derimod hører
eller læser Beretningen og saaledes opfatter Begivenheden
paa anden Haand, kan ikke, som Øienvidnet,
have en umiddelbar Viden af det Oplevede, men maa
nøies med en Tro, en historisk Tro.

Dersom den historiske Tro nu udelukkende hvilede
paa en ubegrændset Tillid til det fortællende Øienvidnes
individuelle Sanddruhed, vilde en saadan Tro umulig
kunne bære nogen videnskabelig Bygning. Historien
vilde paa disse Vilkaar ikke være nogen Videnskab.
% --- p. 315 ---
\setcounter{footnote}{0}%
\opage{315}Men netop den Omstændighed, at her kun er Tale om
Kjendsgjerninger, der ere Naturbetingelser og Menneskelivets
almengjældende Aandslove underlagte, gjør det
muligt, at den historiske Tro kan begrundes, idet Beretningerne
underkastes en Kritik, og deres Troværdighed
prøves.

Troværdigheden af en historisk Begivenhed er tilstrækkelig
begrundet, naar det er gjort indlysende, at
man ikke kan nægte en saadan Begivenhed uden at
nægte al Historie. Sandsynligheden er Troværdighedens
Maalestok: jo sandsynligere en Tildragelse er, desto
troværdigere. Men Sandsynligheden bestemmes igjen ved
de anførte Kjendsgjerningers saavel indre som ydre Sammenhæng. Saaledes kan et Tilfælde, en Handling, et
Karaktertræk, som løsrevet fra Forbindelsen med det
Øvrige, vilde være usandsynligt, ved alsidig, omfattende
Sammenligning blive sandsynligt. Sandsynlighedsgrunde
af prøvet Gyldighed ere tilstrækkelige for den profane
Historie.

Anderledes i den hellige Historie. Her maa Sandsynlighedsgrundene
nedsættes til en underordnet Værdi;
thi det, hvorpaa Alt kommer an, er Underet, og Underet
er ifølge sin Natur det Usandsynlige. Forskjellen
mellem en Undergjerning og en ligefrem historisk Kjendsgjerning
er saa gjennemgribende, at den historiske Tro
forvandler sig og gaaer op i den religiøse. Det er her,
hvad Opfattelsen angaaer, ikke det Ydre, der umiddelbart
bestemmer det Indre, det er omvendt det Indre,
der væsenlig bestemmer det Ydre. Man betragte hvilket
af de evangeliske Undere, man vil, saa viser det sig
strax, at Fænomenet, det Ydre, tilsteder en dobbelt For%
% --- p. 316 ---
\setcounter{footnote}{0}%
\opage{316}klaring. Det kan være en Skuffelse, et Bedrag, et besynderligt
Tilfælde, en overraskende Naturbegivenhed,
og det kan være en Guds Gjerning. Kristus opvækker
de Døde: Jairi Datter (Mtth. 9), Enkens Søn i Nain
(Luk. 7), Lazarus af Bethanien (Joh. 11). Om de to
første Opvækkelser er det nu strax indlysende, at de af
Øienvidnerne kunne forklares paa to Maader: de Ikke-Troende
kunde jo sige, at de saakaldte Døde kun vare
Skindøde. Hvad Lazari Opvækkelse angaaer, da har
den Beskuelighed, hvormed Underet fremtræder, ikke
kunnet forhindre, at en Del iblandt de Tilstedeværende
\emph{have anseet det Hele for et Bedrag.} Det siges jo udtrykkelig,
at Nogle af dem gik hen til Farisæerne og
angave Jesus, og at Farisæerne og Yperstepræsterne netop
paa Grund af denne Angivelse besluttede i Raadet, at
»Folkeforføreren« skulde døe. Betragte vi Opstandelsens
Mirakel\footnote[1]{Jvf. \emph{Om Hindringer og Betingelser for det aandelige Liv i Nutiden.} Efterskrift S. 297—310.}, saa viser sig det Samme. »Der skete
(Mtth. 28) et stort Jordskjælv; thi Herrens Engel nedfoer
af Himmelen, traadte til, og væltede Stenen fra
Indgangen og satte sig der. Men hans Skikkelse var
ligesom Lynet, og hans Klædebon hvidt som Snee. Men
Vægterne skjælvede af Frygt for ham og bleve som
døde«. --- Umiddelbart efter Ordene skulde man nu antage,
at de vagthavende Soldater virkelig havde seet
Engelen i Lynet og saaledes vare blevne Øienvidner til
et virkeligt Under. Men, at dette ikke har været Tilfældet,
fremlyser tydelig deraf, at Vagten lod sig bestikke
af Yperstepræsterne. Havde de Vagthavende
% --- p. 317 ---
\setcounter{footnote}{0}%
\opage{317}anseet det Hele for Andet end en forfærdelig Naturbegivenhed,
da havde de vist ikke havt Mod til at lyve.

Ligger det nu i Sagens Natur, at Opfattelsen af
mirakuløse Begivenheder er betinget af Øienvidnets religiøse
Tro\footnote[1]{Om det logisk Umulige i at opfatte et Mirakel som en anden ligefrem historisk Kjendsgjerning, jvnfr. Grundideernes Logik II. S. 208—210.}, saa kunne Evangelisterne jo heller ikke
indskrænke sig til at berette, hvad de ligefrem have
hørt og seet, men maae fremstille og beskrive, hvad de
\emph{troe at have hørt og seet,} hvad de virkelig have oplevet
i Troen. Den hellige Historie har intet Videns
Grundlag; den er en Historie af det i Troen Oplevede,
en Troens Historie, en Historie af Troende til Troende.

Fra dette Synspunkt vil den i Fortolkerens Grundsætninger
indeholdte Tvetydighed allerede begynde at
vise sig. Den høiærværdige Forfatter fremhæver, at
Evangelisternes Forhold til det Stof, der »udfylder den
evangeliske Beretning«, ikke kan have »samme afgjorte
Fasthed og Bestemthed, som hvor Historieskriveren er
stillet overfor et historisk Stof, hvor Intet gaaer ud
over Grændsen af det Sædvanlige og i sig selv Troværdige«;
han gjør, i god Overensstemmelse med Grundsætningerne
(2) og (3), gjældende, at hvor det historiske
Indhold er af en saa eiendommelig Beskaffenhed, at »det
Hele bevæger sig paa Mysteriets dunkle, hellige Grund«,
der vil »den religiøse Betragtningsmaade ikke kunne
blive uden bestemmende Indflydelse paa den kritiske
Vurdering af Evangeliernes historiske Karakter«. Her
er Tvetydigheden. Det Tvetydige ligger i den Maade,
hvorpaa det Religiøse forbindes med det Videnskabelige.

% --- p. 318 ---
\setcounter{footnote}{0}%
\opage{318}Ved at kalde den religiøse Tro til Hjælp, mener
Fortolkeren at have vundet sikkert Fodfæste mod dem,
der, idet de opfatte Underet som »det ligefrem Naturstridige,
det med Tænkningens som med Naturens
Love Uforenelige«, ville paastaae, »at den evangeliske
Historie ikke kan være foregaaet saaledes, som det,
efter den almindelige Forstaaelse af den evangeliske
Beretning, maatte være Tilfælde«. Her er det nu
strax iøinefaldende, at Forf. i Hovedsagen ikke ret er
uenig med de Modstandere, han af al Magt vil bekæmpe.

De Kritikere, der resolut benægte Underets Mulighed
(Strauss, Weisse o. Fl.), komme med fuld videnskabelig
Følgestrenghed til det velbekjendte Resultat:
»det Guddommelige kan ikke være skeet saaledes; det
saaledes Skete kan ikke være guddommeligt«. Men vor
Fortolker, der har erklæret sig uenig med de moderne Tænkere
angaaende Forudsætningen, er, mærkeligt nok, bleven
enig eller dog næstendels enig med dem i Konsekvensen.
Ved at holde fast paa Adskillelsen af hvad han kalder
»Evangeliernes Udside og det evangeliske Kjernevæsen»
er hans Fortolkning, efter hvad vi ovenfor have seet,
jo i alle væsenlige Punkter en forblommet Stadfæstelse
af den Straussiske Sætning: »det Guddommelige kan
ikke være skeet saaledes«. Forfatteren lover sig noget
Stort af den paaberaabte Adskillelse; ved den »vil det
blive muligt for Apologetiken, seirrigt at afvise de ubeføiede
Slutninger, som den anti-evangeliske Kritik vil
udlede deraf«. Men det vil ikke blive muligt; den antievangeliske
Kritik vil netop seire paa Grund af Adskillelsen;
den vil spørge, og med fuld Ret: »hvad kan det
% --- p. 319 ---
\setcounter{footnote}{0}%
\opage{319}hjælpe, at her antages at være skeet noget Vidunderligt,
naar man dog ikke kan faae at vide, hvad det er,
der er skeet«?

Skal det være mere end en Talemaade, at den hellige
Historie er en Aabenbaringshistorie, da gaaer det
virkelig ikke an at lade det Overnaturlige skjule sig i
et »mystisk Halvmørke« og holde sig som »et tillukket
Mysterium« bagved de sagnagtige Beretninger. At sige
om Evangelisternes Forhold til det givne Stof, Stoffet i
deres Fortællinger, at dette Forhold ikke kan have »samme
afgjorte Fasthed og Bestemthed« som andre Historieskriveres,
er jo rentud at vedgaae, at den evangeliske
Historie, strengt taget, mangler Troværdighed, at Aabenbaringshistorien
er en Historie, der Intet aabenbarer, ---
lucus a non lucendo.

Som Aabenbaringshistorie er den hellige Historie
deri forskjellig fra den profane, at den har \emph{den religiøse}
Tro til sit \emph{oprindelige og selvstændige Erkjendelsesprincip.} Aabenbaring og Tro ere korrelate Begreber.
Som Aabenbaringen er, saaledes er Troen, og som
Troen er, saaledes er Aabenbaringen; begges væsenlig virkende
Enhed er Aanden; men Aandens aabenbarende Udtryk
er Ordet. Lad os da igjen fæste Blikket paa de overnaturlige
Kjendsgjerninger, for at see, hvad heri ligger.

Underet kan ifølge sin Natur umulig være nogen
ligefrem udvortes Erfaringsgjenstand. Det sandselige
Øie seer Lynilden fra Himlen, men Troens Øie seer
Englene ved Graven. Det første umiddelbare Indtryk
er Forfærdelsen; ud af Forfærdelsen fødes Forstaaelsen,
og ud af Forstaaelsen Fremstillingen i Ordet. Det er
Aandens Værk, og som Aanden er, saaledes er Ordet.
% --- p. 320 ---
\setcounter{footnote}{0}%
\opage{320}Er Ordet sagnagtig svævende, uden »Fasthed og Bestemthed«,
saa er Aanden uden Vished, uden Klarhed,
uden Sandhed; men dersom der virkelig er Sandhed,
historisk Sandhed i Aandens Virken, da er der ogsaa
virkelig Sandhed, historisk Sandhed i Ordet.

Grundsætningernes Tvetydighed viser sig nu fra en
ny Side. Den høiærværdige Fortolker føler fuldtvel, at
det blot videnskabelige Princip, der ellers er tilstrækkeligt,
naar det gjælder en Kritik af et Stykke Verdenshistorie,
ingenlunde er tilstrækkeligt, naar det gjælder
en Kritik af den hellige Historie. Saa kalder han
da Troen til Hjælp; men i al Ubestemhed, thi »den religiøse
Betragtningsmaade«, der »ikke kan blive uden
bestemmende Indflydelse paa den kritiske Vurdering af
Evangeliernes historiske Karakter«, er svævende og ubestemt.
Det Videnskabelige er Hovedsagen, Troen kun
en Nødhjælp, der skal læmpes efter Videnskaben.

Ved denne, om end forsigtige, saa dog yderst mislige
Kombineren af Tro og Viden har Fortolkeren beredt
sig selv en stor Forlegenhed. Fra den ene Side
bekæmper han Mythelæren, »den mythiserende Behandling«,
paa det Kraftigste. »En Anstrængelse for at
godtgjøre, at det, som Evangelierne indeslutte og overlevere,
ikke er et ubevidst eller halvbevidst fromt Fantasiværk,
men en Beretning om en virkelig Livshistorie,
virkelige Tildragelser og Gjerninger«, maa, siger han,
»allerede nu agtes ubetimelig«. Hvorfor? Fordi Synsmaaden
er saa urimelig, at det Urimelige indlyser ved
sig selv. Det er jo glædeligt for Troen. Men Glæden
er ikke synderlig stor; thi paa den anden Side bekæmper
han lige saa ivrig den gammelkirkelige Synsmaade,
% --- p. 321 ---
\setcounter{footnote}{0}%
\opage{321}at vi i Evangelisternes Beretninger overalt have »en
virkelig Livshistorie, virkelige Tildragelser og Gjerninger«
for os. Det er en nødvendig Indrømmelse, der
maa gjøres Videnskaben.

Dersom den høiærværdige Forfatter nu blot paa
nogen Maade saae sig i Stand til at forlige de to indbyrdes
stridige Synsmaader med hinanden: da vilde et
saadant Forlig være et talende Bevis for Troens Enhed
med Viden. Men Forliget kommer ikke i Stand. »Paa
samme Tid«, siger han, »som der hæves saa stærke og
saa grundede Klager over den Vilkaarlighed og Uretfærdighed,
for hvilken den hellige Historie er givet til
Pris af den nyeste Kritik --- med al tilsyneladende Strenghed
i den videnskabelige Form, dog i lidet videnskabelig
Aand --- kunne de alvorlige Spørgsmaal ikke skydes
tilbage: om ikke Kritiken er bleven fristet og ægget til
en saa fjendlig Optræden ved gammel, tidt gjentagen og
meget nærgaaende, Udfordring fra Apologetikens Side?
om den apologetiske Opgave kan være rigtig stillet\dots{},
naar den« (Apologotiken) »fremdeles, som hidtil, vil lade   % PRINTED AS IS (Bind 2, p. 321): „Apologotiken“ for Apologetiken (zoom z4)
Bibeltroen træde frem med Magtsprog og Anathema imod
enhver Indrømmelse af Fænomener, der vise hen til en
menneskelig Side af Skriften, fremdeles som hidtil vil
fordre en ubetinget Benægtelse af, hvad enhver sammenlignende
og prøvende Betragtning paanøder os, som
utvivlsom Vished?« I at fremhæve disse Vanskeligheder
har den høiærv. Forf. fuldkommen Ret, og Evangeliekritiken
vil endnu, efterat den længst har forladt hans
Standpunkt, med fornyet Interesse vende tilbage til hans
lærerige Arbeide og altid regne ham det til sand For%
% --- p. 322 ---
\setcounter{footnote}{0}%
\opage{322}tjeneste, at han redelig har modarbeidet Overdrivelserne
og holdt vaagent Øie med Vanskelighederne.

Men at ingen Fortolkning, var det end den lærdeste,
talentfuldeste og skarpsindigste, ud fra det her
valgte Standpunkt er istand til at hæve de paapegede
Vanskeligheder, er klart som Dagen. Hvad kan det
hjælpe at beraabe sig paa en »Adskillelse af Væsenligt
og Uvæsenligt«, naar Adskillelsen, efter hvad vi have
lært af Fortolkerens egne Ord, udtrykkelig viser, at det
Væsenlige er »et tillukket Mysterium«, der ikke kan
antage historisk Form, og Beretningen i sin historiske
Form »den mythiserende Indklædning«, »den udsmykkende
Indfatning«, der gaaer udenom »Kjernen«, det
Uvæsenlige.

Kun paa den Betingelse, at \emph{Troens og dermed}
Aabenbaringens Princip faaer Lov at gjælde for,
hvad det i Sandhed er, \emph{den hellige Histories
eneste afgjørende Erkjendelsesprincip,} kunne
de indviklede Knuder løses. Thi heri ligger: 1) at
Elementerne i den evangeliske Historie ere baade \emph{udvortes}
--- hidrørende fra Alt, hvad der i den ydre
Virkelighed er seet og hørt og erfaret af Øienvidner ---
\emph{og indvortes,} hidrørende fra de Enkeltes særlige Opfattelse
af, hvad de i Troen have hørt og seet og inderlig
oplevet; 2) at de udvortes og indvortes Elementer ere
hos de hellige Vidner saaledes sammensmeltede, sammenarbeidede
og assimilerede; \emph{at ingen Kritik er}
i \emph{Stand til i endelig Forstand at skille dem
fra hverandre;} 3) at det Historiske saavel i Begivenhederne
selv som i de tilsvarende Beretninger hverken
er en blot udvortes ei heller en blot indvortes
% --- p. 323 ---
\setcounter{footnote}{0}%
\opage{323}Virkelighed, men Aandens og Sandhedens Enhed af det
ved sit Indvortes bestemte Udvortes, altsaa \emph{Aandsvirkelighed,
Virkelighed for Aanden.}

Under disse Forudsætninger kunne Vanskelighederne
hæves, thi det Feilsyn, hvoraf de udspringe, kan nøiagtig
paayises. Den hellige Skrift har, hvad vor høiærværdige   % PRINTED AS IS (Bind 2, p. 323): „paayises“ for paavises (zoom z4)
Fortolker særdeles indskærper, ikke blot sin
guddommelige, men ogsaa sin menneskelige Side\footnote[1]{\emph{Om Grundbestemmelser i Aandens Selv,} jvfr. Religionsfilosofi S. 458—486.}. Men
Apologetiken og den negative Kritik ere med Hensyn
hertil midt i al deres Uenighed netop enige i at gaae
ud fra det samme Feilsyn.

Den negative Kritik miskjender Aanden, gaaer bagved
Ordet, prøver de udvortes Kjendsgjerninger og opdager
en saadan Mangfoldighed af Modsigelser og indre
Uoverensstemmelser, at det Hele maa synes utroligt.
Den sætter en selvløs Ide i Aandens Sted og gjør Ordet
til Mythe: her har det Menneskelige kvalt det Guddommelige.

Den orthodoxe Apologetik slaaer sit Dogme fast,
og gjør den evangeliske Historie til et ufeilbart Værk
af Guds ufeilbare hellige Aand. Men til Gunst for sin
Paastand miskjender den Aanden, stirrer paa »det Ydre«,
gaaer bagved Ordet og hænger sig i Bogstaven. Alle
Begivenhederne skulle indtil de mindste Træk være udvortes,
gjengivne med en saa udvortes, fotografisk
Nøiagtighed, at den Hellig-Aand maa have dikteret Alt
indtil Tøddel og Bogstav: her har det Guddommelige
kvalt det Menneskelige.

% --- p. 324 ---
\setcounter{footnote}{0}%
\opage{324}Saa kommer Mæglingsapologetiken. Den vil forlige
det Menneskelige med det Guddommelige, og skynder
sig at begaae den samme Feil som dens Modstandere.
Istedenfor at søge Foreningen af det Guddommelige og
det Menneskelige i Aanden ved Ordet, sønderlemmer
den Ordet, holder sig til Sætninger, overseer Aanden,
det Indvortes, og giver sig paa den allernaiveste Maade
i Kast med det Udvortes. Den søger den egenlige
Kjendsgjerning udenfor Ordet; den titter og speider,
om den dog ikke her eller der kunde skimte, hvad det
Virkelige, det Vidunderlige egenlig er, der staaer bagved
Beretningen. Men al dens Speiden er forgjæves.
Det Virkelige selv, det Vidunderlige, vil ikke lade sig
tilsyne; det er og bliver et ubestemmeligt »Grund-Faktum«,
»et tillukket Mysterium«: Mæglingsapologetiken
opdager Intet.

Af Frygt for »Dualismen« har den mæglende Apologik
hildet sig i den allerfarligste Dualisme, Tvetydighedernes
Dualisme: \emph{een Sandhed i det Videnskabelige,
en anden i det Opbyggelige; eet Evangelium
paa Kathedret, et andet paa Prædikestolen};
\emph{een Tro i Viden, en anden udenfor}
Viden. Skulde det ikke være paa Tide, at vi fik denne
Dualisme opklaret. Tvetydighederne maae bort! Det
er ikke blot et fromt Ønske, et stille Hjertesuk i Troens
Navn; det er Oplysningens Krav, det strenge Bud af den
vrede Fornuft, af Tidens Aand i dens mørke Majestæt.

»\emph{Cardinal, ich habe das Meinige gethan.
Thun Sie das Ihre}«.
\signature{R. \emph{Nielsen}.}

\end{document}
